% Generated draft: requires mathematical and source review.
\documentclass[11pt]{article}
\usepackage[margin=25mm]{geometry}
\usepackage{fontspec}
\setmainfont{Latin Modern Roman}
\usepackage{amsmath,amssymb,mathtools,unicode-math}
\setmathfont{Latin Modern Math}
\usepackage{enumitem,needspace,xurl,hyperref}
\setlength{\parindent}{0pt}
\setlength{\parskip}{5pt}
\setlength{\emergencystretch}{3em}
\newcommand{\N}{\mathbb{N}}
\newcommand{\Z}{\mathbb{Z}}
\newcommand{\Q}{\mathbb{Q}}
\newcommand{\R}{\mathbb{R}}
\newcommand{\C}{\mathbb{C}}
\newcommand{\F}{\mathbb{F}}
\newcommand{\A}{\mathbb{A}}
\newcommand{\PP}{\mathbb{P}}
\newcommand{\E}{\mathbb{E}}
\newcommand{\eps}{\varepsilon}
\newcommand{\dd}{\,\mathrm{d}}
\begin{document}
\section*{UnsolvedMath ID: 800312}
\textbf{Source problem code:} catalogue entry
\subsection{Definitions and mathematical statement}
Let \(S^1=\R/\Z\), and let \(\operatorname{Homeo}(S^1)\) denote its homeomorphism group. For an integer \(r\geq 1\), or \(r=\infty\), let \(\operatorname{Diff}^r(S^1)\) denote the group of bijections whose maps and inverses are \(C^r\). A subgroup \(G\leq\operatorname{Diff}^2(S^1)\) acts minimally if every orbit is dense:
\[
\overline{\{g(x):g\in G\}}=S^1\qquad\text{for every }x\in S^1.
\]
For any subgroup \(H\leq\operatorname{Homeo}(S^1)\), define the normalizer of \(G\) in \(H\) by
\[
N_H(G)=\{h\in H:hGh^{-1}=G\}.
\]
The source asks for a characterization of the minimally acting subgroups \(G\leq\operatorname{Diff}^2(S^1)\) satisfying
\[
N_{H_{\mathrm{ac}}}(G)=N_{H_{\mathrm{diff}}}(G),
\]
where the source describes \(H_{\mathrm{ac}}\) as the group of absolutely continuous homeomorphisms and \(H_{\mathrm{diff}}\) as the group of diffeomorphisms.

These ambient groups require clarification. Absolute continuity of a circle homeomorphism means absolute continuity of its lifts on compact intervals. Requiring absolute continuity of both the map and its inverse gives a group
\[
\{h\in\operatorname{Homeo}(S^1):h\text{ and }h^{-1}\text{ are absolutely continuous}\}.
\]
Requiring it only of the map does not in general give a group. The supplied statement does not explicitly specify the inverse condition. It also does not specify whether \(H_{\mathrm{diff}}\) means \(\operatorname{Diff}^1(S^1)\), \(\operatorname{Diff}^2(S^1)\), or another differentiability class, nor whether orientation-preserving restrictions apply to the acting and ambient groups. Thus the displayed equality records the intended classification question, but a unique precise classification target requires these conventions to be fixed; none is silently imposed here.
\subsection{Short English statement}
Characterize the groups of \(C^2\) circle diffeomorphisms whose every orbit is dense and whose normalizers among absolutely continuous homeomorphisms and among diffeomorphisms coincide. The inverse absolute-continuity convention, ambient differentiability class, and orientation conventions remain unspecified in the supplied source.
\subsection{Sources}
\begin{itemize}
\item Local UnsolvedMath registry, permanent ID \(800312\), source problem code \texttt{catalogue entry}, title ``Group Actions on 1-Manifolds --- Question 29.'' The record attributes the question to Navas, with proposed year 2017, and labels it open; this status is supplied metadata, not an independently verified assessment.
\item Navas, \emph{Group Actions on 1-Manifolds: A List of Very Concrete Open Questions} (2017), source item recorded as ``Question 29 (\#28).'' Supplied source URL: \url{https://arxiv.org/abs/1712.06462}. The record reports access on 2026-10-16 and extraction by \texttt{arxiv\_latex}; these are supplied provenance statements, not a claim of independent access or verification.
\end{itemize}
\end{document}
